\section{Discussion}
\chead{\thesection. Discussion}

\subsubsection*{Sparsity as a Computational Enabler}

Only a tiny fraction of all possible protein pairs are preserved after
pairwise alignment and the subsequent E-value filtering: for the full \texttt{ref\_prot} dataset (82.1 million
proteins, 1.76 billion edges), the fraction of observed edges relative to a
complete graph is approximately \(5.2\text{e}-7\). This extreme sparsity makes
graph‑based clustering feasible on a single node, as it reduces both memory and
time complexity from \(\mathcal{O}(n^2)\) and \(\mathcal{O}(n^3)\) to
\(\mathcal{O}(n k_0)\) and \(\mathcal{O}(n k_0^2)\), respectively, where \(k_0
\approx 21\) is the average number of non‑trivial entries per column of the
adjacency matrix. The complexity analysis in the Background
(Observe Fig.~\ref{fig:complexity}) quantifies this advantage and is directly reflected
in the experimental runtimes: MCL completes within 2–3 days for the full
dataset using \(\sim 440\) GiB of memory, while a dense approach would
require petabytes of memory and years of computation.

\subsubsection*{Qualitative Characteristics of the Data and Clusters}

Frequency diagrams of amino acid compositions shown in Fig.~\ref{fig:aa_distr}
and sequence lengths of Fig.~\ref{fig:len_freq} provide a global signature of
each selected proteome but are not sufficient for clustering: they describe global
biases, not the pairwise relationships that determine cluster membership.
Additionally, from Fig.~\ref{fig:aa_distr}, the observed hydrophobic‑to‑charged
residue ratios range from 1.14 \textit{Danio rerio} (zebrafish) to 1.46
\textit{Carnobacterium maltaromaticum} (bacterium). To conveniently
discriminate between hydrophobic and charged residues, the reader can also view
the amino acid lookup table in Appendix~\ref{appendix:aminoAcidLookup}. 

Recent approaches have used learned embeddings from protein language models,
including TAPE~\cite{rao2019transformer}, ESM-2~\cite{lin2022esm2}, and
ProGen2~\cite{nijkamp2022progen2}, to represent sequences in Euclidean space.
While these methods achieve strong performance on some tasks, they trade off
interpretability and introduce large memory overhead. In contrast, the
graph‑based representation used here derives similarity directly from alignment
scores, preserving biological meaning without black‑box abstraction.

With this graph‑based representation, the resulting clusters can be directly
interpreted by examining their taxonomic composition and structural coherence.
The three largest clusters each consist predominantly of proteins from the same
organism \textit{Austropuccinia psidii MF-1}, (rust fungus) with TaxID 01389203
a strain from a species that causes myrtle rust in various plants.
\textit{Austropuccinia psidii} (myrtle rust) is a pathogen of Myrtaceae plants
\cite{psidii}. This can be observed in the bar plot of
Fig.~\ref{fig:partition_composition} (j), showing the top-25 clusters ranked
from largest to smallest. Hence, this suggests that this organism possesses
unusually large and diverse protein families, each large enough to form a top
cluster. Other organisms appear primarily in smaller clusters, indicating
either lower representation in the dataset or less internal sequence diversity.
Purity values are high across the top‑25 clusters, indicating that MCL tends to
group proteins that match the biological profile of a limited set of dominant
organisms, especially in the largest clusters.

\subsubsection*{Algorithmic Behaviour: Granularity and Differences Between MCL and HipMCL}

Higher inflation consistently produces more clusters than lower inflation as
shown in Table~\ref{tab:experimentsPalladium} and, for datasets up to 35
million proteins, also reduces runtime. This matches MCL's design: inflation
strengthens dominant connections, accelerates convergence, and yields finer
granularity. The cluster size distribution (Fig.~\ref{fig:cluster_distr})
partially supports this premise: the distribution moves towards smaller
clusters when inflation is set to 2.5 for both algorithms.

Although MCL and HipMCL share the same core algorithm, HipMCL consistently
produces slightly more clusters at the same inflation value, likely due to
differences in pruning strategies or numerical precision. However, its memory
overhead (2–4 times higher than MCL) and convergence issues make it impractical
for single‑node runs on datasets larger than 3 million sequences. This does not
diminish HipMCL's value; it was designed for distributed environments, where it
excels.

\subsubsection*{Reproducibility and Random Walks}

Both MCL and HipMCL rely on random walks that simulate flow in the graph, but
their implementations differ in pruning strategies, numerical precision, and
convergence criteria. Table~\ref{tab:experimentsPalladium} shows that for the
same inflation, direct quantitative comparison of individual clusters between
MCL and HipMCL is inherently ambiguous.

However, this does not compromise the reproducibility of the present work. The
reported metrics --- runtime, peak memory, number of clusters, cluster size
distributions, and taxonomic purity --- are all stable across runs. Moreover,
the exact clusters produced by each algorithm can be reproduced by fixing the
input graph and using the same software versions (commit digests) and
parameters. The workflow is implemented in \texttt{bioscflow}, a Python
codebase that orchestrates GeneCAST, DIAMOND, MCL, and HipMCL (observe
Fig.~\ref{fig:workhigh}), manages file paths and intermediate data, and
systematically records execution logs and runtime metrics in JSON session
reports which is particularly useful when producing multiple experiments. The
observed differences in cluster counts and granularity between MCL and HipMCL
are therefore meaningful as algorithmic characteristics, even if individual
cluster assignments are not
directly comparable.

\subsubsection*{Limitations}

Three main limitations should be noted. First, HipMCL was tested only in
single‑node mode; its poor performance here does not reflect its capabilities
on large clusters. Second, biological validation of individual clusters was
outside the scope of this work. Third, the reference proteome dataset, while
large, is biased toward well‑studied organisms, and the chosen BLASTp
thresholds may not be optimal for every application.

Another major limitation concerns the stochastic nature of the random walks
underlying both MCL and HipMCL. Because the algorithms rely on simulated flow,
runtime measurements may exhibit non‑deterministic fluctuations. In principle,
a larger input could occasionally complete faster than a smaller input in a
single run, potentially obscuring true scaling behavior. All reported runtimes
are based on single executions; therefore, the observed trends should be
interpreted with caution. Systematic averaging over multiple runs with
different random seeds would be required to obtain statistically reliable scaling
estimates.

\subsubsection*{Design Priorities Revisited}

The workflow meets the three priorities set in the Introduction, namely
\textit{Scalability}: MCL processed 82 million sequences on a single node.
\textit{Reproducibility}: the entire pipeline is implemented in
\texttt{bioscflow} with versioned tools and commits. \textit{Interpretability}:
graph edges are derived directly from alignment scores, avoiding black‑box
embeddings.
